\documentclass[10pt,a4paper]{article}

\usepackage[margin=1.05in]{geometry}
\usepackage{setspace}
\usepackage{amsmath,amssymb,amsthm}
\usepackage{graphicx}
\usepackage{enumitem}
\usepackage{titlesec}
\usepackage[expansion=false]{microtype}
\usepackage{fancyhdr}
\usepackage{hyperref}
\hypersetup{colorlinks=true,linkcolor=black,urlcolor=black,citecolor=black}
\graphicspath{{/home/workdir/artifacts/}}

\titleformat{\section}{\normalsize\bfseries}{}{0em}{}
\titleformat{\subsection}{\small\bfseries}{}{0em}{}
\titlespacing*{\section}{0pt}{1.35em}{0.5em}
\titlespacing*{\subsection}{0pt}{0.9em}{0.3em}

\theoremstyle{plain}
\newtheorem{proposition}{Proposition}
\theoremstyle{definition}
\newtheorem{definition}{Definition}
\theoremstyle{remark}
\newtheorem{observation}{Observation}

\setstretch{1.14}
\setlist[itemize]{leftmargin=1.15em,itemsep=0.1em,topsep=0.2em}

\pagestyle{fancy}
\fancyhf{}
\fancyhead[L]{\small\itshape Elements of Position}
\fancyhead[R]{\small\itshape On a Granted Power}
\fancyfoot[C]{\small\thepage}
\renewcommand{\headrulewidth}{0pt}

\begin{document}

\begin{center}
\vspace*{1.6em}
{\large Elements of Position}\\[1.4em]
{\LARGE\bfseries On a Granted Power}\\[0.5em]
{\small The next pair revises the tilt. Equal weight is the only rate that stays with the list.}\\[1.5em]
{\normalsize Justin Erholtz}\\[0.3em]
{\small 20 September 2026}
\end{center}

\vspace{1.2em}

\noindent
\textit{A note.}
Part~I wrote the pair on the hit count \(m\): inner \(1/m\), outer \(m\), product \(1\).
Part~II wrote the kite factor \(1-1/m^2\) and the increment of that list when hits are kept or sorted.
This part grants a power \(s\) on the same inner and records what the next hit does to a zero of the resulting balance.
Claims of two kinds stay apart.
An identity is checked against the sum.
A path of a zero as \(N\) grows is a table.
The operators do not contain the power.
The pair contains the inner.

\vspace{0.6em}
\noindent
\textit{A note on names used here.}
Inner, outer, lock, hit, pair: as in Part~I.
Power \(s\): a chosen complex number.
Balance \(B_N(s)\): defined below.
Tilt: a zero of \(B_N\) with real part not equal to \(1/2\).
Landing: that real part becoming \(1/2\) at a later hit, and staying there while the tracker holds.

% ======================================================================
\section{What is granted}
% ======================================================================

The inner at hit \(m\) is \(1/m\).
Write \(m^{-s}\) for a chosen \(s\).
At \(s=2\) this square is already the inner square in the kite factor \(1-1/m^2\).
At every other \(s\) the power is extra.

Two writings of one inner: \(m^{-s}\) and \(m^{-(1-s)}\).
Swapping inner with outer is the map \(s\mapsto 1-s\).
That is the inversion of Part~I, read on the power.

\begin{definition}[Balance]
\[
B_N(s)=\sum_{m=1}^{N}\bigl(m^{-s}-m^{-(1-s)}\bigr).
\]
\end{definition}

\begin{proposition}
For every \(N\ge 1\) and every \(s\),
\[
B_N(1-s)=-B_N(s).
\]
\end{proposition}

\begin{proposition}
\[
B_N\bigl(\tfrac12\bigr)=0.
\]
Each term vanishes separately: \(m^{-1/2}-m^{-1/2}=0\).
\end{proposition}

The lock term is identically zero for every \(s\).
The first split is the first nonzero term.
That term alone vanishes only at \(s=1/2\).
A zero with real part not \(1/2\) therefore uses more than one post-lock hit.

% ======================================================================
\section{The increment}
% ======================================================================

Fix a tilt \(s_N\) of the first \(N\) pairs:
\[
B_N(s_N)=0,\qquad \mathrm{Re}(s_N)\ne\tfrac12.
\]

Hit \(N+1\) adds one pair, lengths \(N+1\) and \(1/(N+1)\).
Its term at the old tilt is
\[
\tau_{N+1}(s)=(N+1)^{-s}-(N+1)^{-(1-s)}.
\]

\begin{proposition}[Increment]
\[
B_{N+1}(s_N)=\tau_{N+1}(s_N).
\]
\end{proposition}

\begin{proof}
\[
B_{N+1}(s)=B_N(s)+\tau_{N+1}(s).
\]
At \(s=s_N\) the first summand is zero.
\end{proof}

Checked on the recorded path called L below, at \(N=128,140,150,160,163\): the two sides agree to the digits printed in that run.

The old tilt is therefore not a zero of \(B_{N+1}\).
The leftover is exactly the new pair’s term.
A new tilt \(s_{N+1}\) is any number with \(B_{N+1}(s_{N+1})=0\).
To first order,
\[
s_{N+1}-s_N\approx -\frac{\tau_{N+1}(s_N)}{B_N'(s_N)}.
\]
On the recorded paths that estimate tracks the change in the imaginary part to about \(3\times 10^{-4}\) per step while the step is small.
It is only rough on the real part, and it fails at a landing step, where the real part jumps to \(1/2\) in one hit.

After a landing the new-pair term is still not zero.
It no longer changes the real part.
It changes the imaginary part.

% ======================================================================
\section{What the next pair does: two recorded paths}
% ======================================================================

Path L. Started at \(N=128\), \(s=0.803+19.398i\).
Each step uses the previous zero as the Newton start.

\begin{center}
\begin{tabular}{rccc}
\(N\) & \(\mathrm{Re}(s)\) & \(\mathrm{Im}(s)\) & \(\lvert\mathrm{Re}(s)-1/2\rvert\) \\
\hline
128 & 0.803 & 19.398 & 0.303 \\
160 & 0.610 & 18.496 & 0.110 \\
164 & 0.500 & 18.424 & 0 \\
320 & 0.500 & 16.824 & 0
\end{tabular}
\end{center}

Landing at \(N=164\).
Every recorded step before that landing decreased the distance to \(1/2\).

Path G. Started at \(N=96\), \(s=0.655+17.256i\).
Same continuation.

\begin{center}
\begin{tabular}{rcccc}
step & \(\mathrm{Re}(s)\) & \(\Delta\mathrm{Re}\) & \(\lvert\tau\rvert\) & direction \\
\hline
\(96\to 97\) & 0.655 & \(+0.017\) & 0.175 & away \\
\(110\to 111\) & 0.752 & \(+0.002\) & 0.284 & away \\
\(120\to 121\) & 0.756 & \(-0.001\) & 0.284 & in \\
\(160\to 161\) & 0.633 & \(-0.006\) & 0.117 & in \\
\(172\to 173\) & 0.527 & \(-0.027\) & 0.032 & in, lands
\end{tabular}
\end{center}

Twenty steps away from \(1/2\), then fifty-seven in.
The new-pair term grew on the away stretch and shrank on the in stretch.
The first-order sign of \(\Delta\mathrm{Re}\) matched the away stretch.
It lagged the actual turn by a few hits.
Landing at \(N=173\).

\begin{figure}[htbp]
\centering
\includegraphics[width=0.92\textwidth]{GP1_paths_sigma.png}
\caption{Sheet GP1. Real part of paths G and L against the hit count. The rule is \(\mathrm{Re}(s)=1/2\).}
\end{figure}

\begin{figure}[htbp]
\centering
\includegraphics[width=0.92\textwidth]{GP2_plane.png}
\caption{Sheet GP2. The same paths in the plane of the granted power.}
\end{figure}

The increment identity does not have one sign.
The next pair can raise the real part or lower it.
It cannot leave the old tilt a zero.
The new length is fixed.

% ======================================================================
\section{Table of continued families}
% ======================================================================

Nineteen families were started from off-line zeros found at \(N=32\), \(96\), \(128\), and \(192\), in a search window \(\mathrm{Re}(s)\in[-0.4,2.4]\), \(t\in(0,42]\), and continued hit by hit.

\begin{center}
\begin{tabular}{lccc}
born at & families & landing \(N\) & later status \\
\hline
32 & 4 & 37, 39, 47, 56 & on the line through 240 \\
96 & 3 & 97, 97, 173 & on the line; one tracker jump \\
128 & 8 & 133--179 & on the line through 450 \\
192 & 4 & 202--303 & on the line; one tracker jump
\end{tabular}
\end{center}

Nineteen started.
Nineteen landed at a finite later \(N\).
Two trackers jumped to a distant imaginary part after landing and were dropped.
No continued family returned to \(\mathrm{Re}(s)\ne 1/2\) while the residual stayed below \(10^{-6}\).

Two families born at different \(t\) were observed to occupy the same on-line point at large \(N\).

This is a table of families that were started.
It is not a census of every zero of every \(B_N\).
A min-over-grid distance to \(\mathrm{Re}(s)=1/2\) is not monotone in \(N\), because new off-line zeros appear as \(N\) grows and because a family already in hand can move away before it lands (path G).

\begin{observation}[A wider census, checked against the argument principle]
\label{obs:widecensus}
A grid search is only ever a sample; it cannot certify that nothing was missed.
The argument principle can: for an entire function, the winding of its values around a rectangle's boundary counts the zeros inside \emph{exactly}, independent of any grid.
At the same seven birth points, \(N=32,64,96,128,160,192,256\), and the same window, \(\sigma\in[0.51,2.4]\), \(t\in[0.5,90]\) --- more than double the earlier search in \(t\), against four birth points, not seven --- the winding of \(B_N\) around that rectangle gives the exact off-line count at each \(N\): \(14,27,29,28,33,32,34\), a total of \(197\).
A first search of this window, by grid and Newton refinement at a fixed resolution, found only \(109\): roughly half.
A second search, refining not by a finer fixed grid but by recursive subdivision --- halve any sub-rectangle whose winding is not \(0\) or \(1\), refine by Newton once it is \(1\), and accept the root only if it lands back inside the sub-rectangle that isolated it --- found \(192\) of the \(197\), exact at two of the seven birth points and short by one at the other five.
The five short cases trace to a sub-rectangle boundary passing close enough to a zero to make its winding briefly ambiguous at the resolution tried; deeper subdivision would likely close the rest, and does not change what follows.
All \(192\) were continued, hit by hit, to \(N=2000\).
Every one landed.
None was observed to leave \(\mathrm{Re}(s)=1/2\) again while its residual stayed below \(10^{-6}\).
Landing \(N\) ranged from \(33\) to \(393\), median \(166\).
\end{observation}

This does not change what the original table already stood on --- landing is still reported from a table, not proved --- but the table is now checked, not merely sampled: \(197\) is not a guess at how many off-line zeros the window holds, it is what the window holds, by the same principle that makes a contour integral count something a grid can only estimate.

\begin{proposition}[Exponential type of the line reduction]
\label{prop:linetype}
On \(\mathrm{Re}(s)=1/2\), write \(s=\tfrac12+it\).
Then \(B_N(\tfrac12+it)=-2i\,f_N(t)\), where
\[
f_N(t)=\sum_{m=2}^N m^{-1/2}\sin(t\ln m).
\]
\(f_N\) extends to an entire function of the complex variable \(t\), of exponential type at most \(\ln N\).
\end{proposition}

\begin{proof}
\(m^{-1/2-it}-m^{-1/2+it}=m^{-1/2}\bigl(e^{-it\ln m}-e^{it\ln m}\bigr)=-2i\,m^{-1/2}\sin(t\ln m)\); sum over \(m=2,\ldots,N\) (the lock term \(m=1\) contributes \(\sin(0)=0\)).
For complex \(t=x+iy\), \(|\sin(t\ln m)|\le\tfrac12 e^{|y|\ln m}\le\tfrac12 e^{|y|\ln N}\) for every \(m\le N\).
Summing,
\[
|f_N(x+iy)|\le\tfrac12\,e^{|y|\ln N}\sum_{m=2}^N m^{-1/2},
\]
a bound growing at rate \(\ln N\) in \(|y|\) and independent of \(x\); this is exactly the definition of exponential type \(\le\ln N\).
\end{proof}

By the density theorem for entire functions of exponential type bounded on the real line \cite{boas1954}, an entire function of type \(\tau\) has, in any interval of length \(L\), at most \(\tau L/\pi+o(L)\) real zeros --- the same bound \(\sin(\tau t)\) meets with equality.
Proposition~\ref{prop:linetype} therefore forces the density of on-line zeros of \(B_N\) to be at most \((\ln N)/\pi\) per unit length of \(t\), for every \(N\).
This is not the plausibility argument of an earlier draft of this note; it is a consequence of \(f_N\) itself, checked once and standing for every \(N\).

\begin{observation}[The bound holds, and is not tight]
\label{obs:crowding}
Measured directly --- counting sign changes of \(f_N\) over long windows of \(t\), averaged over several windows to damp sampling noise --- the on-line zero density grows with \(N\), stays under Proposition~\ref{prop:linetype}'s ceiling at every \(N\) tried, and does not approach it:
\begin{center}
\begin{tabular}{rccc}
\(N\) & measured density & ceiling \(\ln(N)/\pi\) & ratio \\
\hline
16 & 0.61 & 0.88 & 0.69 \\
64 & 1.01 & 1.32 & 0.76 \\
256 & 1.18 & 1.77 & 0.67 \\
1024 & 1.20 & 2.21 & 0.54 \\
4096 & 1.29 & 2.65 & 0.49 \\
\end{tabular}
\end{center}
The ratio falls as \(N\) grows: the true density grows more slowly than the ceiling, by a factor this note does not derive.
What the ceiling already explains is the wander seen in continuing a landed family well past its landing \(N\): one family, born at \(N=64\), landed at \(N=65\) near \(t\approx30.7\); by \(N=100\) it sat near \(t\approx45.5\) and stayed there, to within about \(0.2\), through \(N=800\), \(\mathrm{Re}(s)\) at \(1/2\) to machine precision throughout, before a Newton tracker with a bounded iteration budget lost the thread.
A window of \(t\) that holds one on-line zero at \(N=65\) holds several by \(N=800\), by a growth Proposition~\ref{prop:linetype} bounds and the table above measures; at that density a step taken from one on-line zero can land on a neighbor instead of re-finding the same one, and asking which of them is still ``the same'' family is a question the operators do not answer.
This, not a return off the line, is why most trackers in Observation~\ref{obs:widecensus} eventually lost lock, always long after that family had already landed --- and why raising the tracker's own iteration budget recovers some of them, since the difficulty is telling neighbors apart, not losing the line.
\end{observation}

% ======================================================================
\section{Where the list does not halt}
% ======================================================================

Infinity in this construction is already written. It has four readings. They are not four objects.

\begin{proposition}
\[
B_N(1)=\sum_{m=1}^{N}\frac1m - N,\qquad
B_N(0)=N-\sum_{m=1}^{N}\frac1m=-B_N(1).
\]
\end{proposition}

At \(s=1\) one writing is the inner and the other is the constant \(1\).
At \(s=0\) those two writings have swapped.
The lock term is still zero.
The first split contributes \(-1/2\) at \(s=1\) and \(+1/2\) at \(s=0\).
As \(N\) grows both quantities leave every bound: at \(N=127\), \(\lvert B_N(1)\rvert\approx 121.57\); at \(N=400\), \(\approx 393.43\).
The infinite list does not produce a finite balance at these two powers.
They sit at the ends of the segment that \(1/2\) bisects.

The hit count itself has no last term.
Inner \(1/m\to 0\), outer \(m\to\infty\).
Inversion \(Q=P/|P|^2\) has its pole at the origin.
The walk never occupies that point: \(\lvert P_m\rvert=m\ge 1\).
As \(m\to\infty\), \(P\) leaves every disk and \(Q\) approaches the origin.
That pole is the same infinity as the hit count, read on the ray.

Exact closure of Part~I is the kick \(h=\tan(\pi/n)\).
At \(n=2\) the formula is undefined.
A large finite kick does not write that half-turn.
On trials \(h=2,5,10,20,50,100\), the first sign change of the second coordinate from the start \((1,0)\) occurs at step \(3\).
The heading's second coordinate at that hit approaches \(-1\) and does not arrive at a finite kick.
That pole is infinity of the kick, not of the pair.

% ======================================================================
\section{Which weight can stay a balance}
% ======================================================================

The two writings at hit \(m\) have sizes \(m^{-\sigma}\) and \(m^{\sigma-1}\).
Those sizes are equal if and only if \(\sigma=1/2\).
The larger size is \(m^{-\alpha(\sigma)}\) with
\[
\alpha(\sigma)=\max(\sigma,\,1-\sigma).
\]
Always \(\alpha(\sigma)\ge 1/2\), with equality only at \(\sigma=1/2\).

At a frozen \(s\) the next pair therefore arrives with size about \(N^{-\alpha(\sigma)}\).
On the equal-weight line that is \(N^{-1/2}\).
Off the line it is strictly larger.

Recorded \(\lvert B_N\rvert\) at \(t=18\) through \(N=2048\):
on \(\sigma=1/2\), \(\lvert B_N\rvert/\sqrt{N}\) stays of order \(10^{-2}\) and does not trend up;
on \(\sigma=0.8\), \(\lvert B_N\rvert/N^{0.8}\) stays of order \(5\times 10^{-2}\), so \(\lvert B_N\rvert\) grows like \(N^{0.8}\).
The path at \(\sigma=0.2\) is the same numbers, by reflection.

The infinite list can keep an oscillating partial sum at the on-line rate.
Off the line the partial sum follows the faster writing.
A frozen unequal weight is not a balance of the whole list.
A tilt has to move.
This does not force every moving tilt to arrive at a finite \(N\).
It does force the comparison: a weight that stays off \(1/2\) pays a strictly larger term at every later hit.

\begin{figure}[htbp]
\centering
\includegraphics[width=0.92\textwidth]{GP3_growth.png}
\caption{Sheet GP3. Frozen \(s\), \(t=18\). On-line scaled by \(\sqrt{N}\); off-line scaled by \(N^{0.8}\).}
\end{figure}

% ======================================================================
\section{What stands}
% ======================================================================

Forced, once the power is granted:
the reflection \(B_N(1-s)=-B_N(s)\);
the termwise zero at \(s=1/2\);
the increment \(B_{N+1}(s_N)=\tau_{N+1}(s_N)\);
the lock term identically zero;
the two closed forms \(B_N(1)=\sum 1/m-N\) and \(B_N(0)=N-\sum 1/m\);
\(\alpha(\sigma)=\max(\sigma,1-\sigma)\ge 1/2\), equality only at \(\sigma=1/2\);
the exponential type of the line reduction \(f_N\) (Proposition~\ref{prop:linetype}), and the ceiling \((\ln N)/\pi\) it forces on on-line zero density.

Reported from a table:
the wider census of Observation~\ref{obs:widecensus} --- \(197\) off-line zeros in the search window, exactly, by the argument principle; \(192\) of them found and continued; all \(192\) landed by \(N=393\);
the away-then-in stretch on path G;
the first-order estimate’s range of usefulness;
\(\lvert B_N\rvert\) at frozen \(s\), on and off the line, through \(N=2048\);
the measured on-line density against its proven ceiling (Observation~\ref{obs:crowding}), well under at every \(N\) tried and not shown to approach it.

Not granted:
a landing of every zero of every \(B_N\);
a value of \(s\) treated as written by the shear;
a motion of the kite, which does not depend on \(s\);
a fixed identity for an on-line family once the zeros crowding its window (Observation~\ref{obs:crowding}) make ``the same zero'' a question the operators do not answer;
the true growth rate of that crowding, only a proven ceiling on it and a measurement sitting under that ceiling.

The pair wrote the inner.
The next hit wrote a fixed new inner.
The increment is that new inner, raised to the two writings, and nothing else.
Where is \(1\)?
It is still the lock.
The lock does not enter the balance.
Everything after is the same pair, with a power placed on it, and with one more count.

\newpage
\begin{thebibliography}{9}

\bibitem{boas1954}
R.\,P.\ Boas, Jr., \emph{Entire Functions}, Academic Press, 1954.

\end{thebibliography}

\end{document}
